\chapter{What this is}

This chapter states the construction, places it beside the crystallographic object it resembles, and marks the boundary of what the work can claim about it.

\section{The construction}

Tile a published unit cell, carry every site as an exact integer, and score how well the arrangement agrees with itself shifted by each candidate offset along one axis. The offset where agreement peaks is the cell edge.

That is an autocorrelation. For one lag along one axis, the score is the share of \emph{occupied} sites whose neighbor at that distance carries the same value. The candidate lags are swept and the score is read at each.

Empty positions are left out, and the measurement depends on that exclusion. A structure built from atom sites is mostly empty space, and counting the empty positions would report how empty the structure is, which every structure agrees on.

\section{The object it resembles}

The autocorrelation of electron density in a crystal is the Patterson function~\cite{src:Patterson-1934}, whose peaks sit at the vectors between atoms. The construction here is a one-axis, element-labeled relative of it: it is read on sites and not on a density, two sites agree only when they carry the same element, and the cell edge is the lattice vector at which every site agrees.

No priority is claimed for the construction. It is reached from a general question, how far an arrangement sits from the most disordered arrangement of its own parts, with no crystal in mind. Arriving from there at an object a field already uses is evidence that the general question is pointed at something real. A full placement of this construction in the crystallographic literature is not made here.

\section{Why the crystal is here as a control}

The interest is not the cell edge. Cell edges are published, refined to more decimal places than anything here reaches, and recovering one adds nothing to crystallography.

The interest is that the answer is known in advance by somebody else. Apart from the protein structures, every other control in this work is a memoryless process, and a memoryless control shows only that an instrument does not invent structure. It cannot show that the instrument finds structure that is there. Those are different claims, and an instrument has to earn the second.

A crystal earns it. The periodicity is not inferred from the measurement; it is a fact about the mineral, and the Crystallography Open Database~\cite{src:Crystallography-Open-Database} publishes the cell for every entry it holds. The number the detector has to hit is written down by people with no interest in whether it works.
